\section{Fourth-order confidence geometry}
\label{sec:compensated}
The positive smoothing kernel in Section~\ref{sec:clouds} uses only a
third-derivative bound. The class of Definition~\ref{def:prior} has six
bounded derivatives. Exploiting those derivatives reduces both the number
of geometric launches and the localization precision required by the
same certificate. No additional observation or regularity assumption is
introduced. The positive-kernel construction remains available under its
weaker local smoothness hypothesis.

Let $\varphi$ be the even kernel in Section~\ref{sec:clouds}, and define
\begin{equation}\label{eq:comp-kernel}
 K_h=\tfrac43\varphi_h-\tfrac13\varphi_{2h},\qquad 0<h\le\tfrac12.
\end{equation}
The periodic rescalings have support in an interval of radius $2h$ before
periodization. This is a signed kernel. Its zeroth moment is one and its
moments of orders one, two and three are zero. Indeed the odd moments
vanish by symmetry, while the second moments of $\varphi_h$ and
$\varphi_{2h}$ are in the ratio $1:4$.

\begin{lemma}[Compensated support approximation]\label{lem:comp-smoothing}
Let $p$ be a periodic support function with $\|p\|_{C^6}\le B_6$, and
suppose $\|p_P-p\|_\infty\le e$. With
$C_\varphi=1+\|\varphi'\|_1+\|\varphi''\|_1$ and
$\widetilde p=K_h*p_P$, one has
\begin{equation}\label{eq:comp-bound}
 \|\widetilde p-p\|_{C^2}
       \le 2C_\varphi e h^{-2}+B_6h^4.
\end{equation}
Here the $C^j$ norm is the maximum of the sup norms of derivatives
through order $j$. If the right side is less than $1/(4\kappa_+)$
and $p+p''\ge1/\kappa_+$, then
$\widetilde p+\widetilde p''>1/(2\kappa_+)$. In particular the
signed correction still defines a regular strictly convex closed curve.
\end{lemma}
\begin{proof}
For $0\le j\le2$,
\[
 \|K_h^{(j)}\|_1\le
 \left(\tfrac43+\tfrac13\,2^{-j}\right)
       h^{-j}\|\varphi^{(j)}\|_1\le 2C_\varphi h^{-2}.
\]
This proves the noise term without differentiating the polygonal support
function. Apply Taylor's formula of degree three to $p^{(j)}$.
The moment identities cancel every nonconstant polynomial term. The
absolute fourth moment is bounded by
\[
 \int |t|^4|K_h(t)|\,\dd t
 \le \left(\tfrac43+\tfrac{16}3\right)h^4
                            \int u^4\varphi(u)\,\dd u
 \le \tfrac{20}3h^4.
\]
The remainder is consequently bounded by
$(20/72)B_6h^4\le B_6h^4$, including $j=2$.
The periodic statement follows by using the periodic lift of $p$ in
this local Taylor calculation. Finally, the errors of $p$ and $p''$
add to at most twice the displayed $C^2$ bound. The positive lower bound
for the radius of curvature proves strict convexity and regularity.
Convexity here follows from this inequality, not from averaging convex
bodies with positive coefficients.
\end{proof}

\begin{theorem}[Reduced geometric launch budget]\label{thm:comp-probe}
Use the full $C^{6,\beta}$ class of Definition~\ref{def:prior}. The
outputs and conditional confidence guarantee of
Theorem~\ref{thm:probe-main} can instead be obtained with
\begin{equation}\label{eq:comp-probe}
 C\nu^{-3/4}\log\frac{C}{\nu\zeta}
\end{equation}
auxiliary launch attempts and impact-position error tolerance at most
$c\nu^{3/2}$. The constants depend only on the same numerical priors.
Neither an analytic continuation constant nor a boundary-query oracle
is used. Deterministic convolution and quadrature errors are included
in the geometric error budget; their arithmetic work is not included
in the launch count.
\end{theorem}
\begin{proof}
Translate the cloud by one of its measured seeds. One may take
$B_6=2+M_6+2D_0$ for the true support in this frame. For all sufficiently
small $\nu$, use
\begin{equation}\label{eq:comp-choices}
 \begin{gathered}
 h=(\nu/(8B_6))^{1/4},\qquad
 E_\nu=\frac{\nu h^2}{128C_\varphi},\\
 q=\min\{d_0/80,1/2,\pi/(2\kappa_+),
                         \sqrt{E_\nu/(4\kappa_+)}\},\\
 \varepsilon_x\le\min\{E_\nu/8,q/20,d_0/200\}.
 \end{gathered}
\end{equation}
The remaining bounded range uses the smaller fixed thresholds. The
confidence hull of Lemma~\ref{lem:cloud} has error
$e_q\le E_\nu/4$. Equation~\eqref{eq:comp-bound} gives
\[
 2C_\varphi e_qh^{-2}+B_6h^4
                 \le\nu/256+\nu/8.
\]
Allow an additional $\nu/16$ for numerical convolution and integration.
The total is below $\nu/4$. For small $\nu$ the preceding lemma
therefore supplies a strictly convex support function, and the intrinsic
coordinate, normal-root, fingerprint and clearance arguments apply with
the same errors. Distance-to-solid enclosures continue to use the original
hulls and \eqref{eq:distance-cloud}, rather than an unjustified
monotonicity of the signed reconstruction.

In the asymptotic range $E_\nu=c_1\nu^{3/2}$ and
$q=c_2\nu^{3/4}$. Substitution into the exact first-hit coverage budget
\eqref{eq:sample-budget} proves \eqref{eq:comp-probe}. These are
attempts, including solid initial positions and no-impact trajectories.
The confidence argument is unchanged after conditioning on any preceding
history: the square and bodies are fixed and the auxiliary launches are
fresh. Every new smoothness bound used here is already supplied by
Definition~\ref{def:prior}. No differentiability of the noisy hull is
assumed.
\end{proof}

\begin{corollary}[Reduced total launch cost]\label{cor:comp-work}
In Theorem~\ref{thm:main}, replace each positive-kernel cloud call by
Theorem~\ref{thm:comp-probe}, keeping the same raw proposals, fresh
validation, error spending and schedule \eqref{eq:schedule}.
The anytime certificate, the stopping tail \eqref{eq:tail-main} and
the certified reconstruction error are unchanged. The cumulative launch
bound improves from \eqref{eq:cost-main} to
\begin{equation}\label{eq:comp-work}
 C(m+1)^{19/4}\log\frac{C(m+1)}{\delta},
\end{equation}
with the same separate rejection-tail allowance. Impact localization
need only satisfy $\varepsilon_{x,k}\le c\nu_k^{3/2}$.
The expected total number of launches up to $T_\xi$ is finite whenever
the spending exponent satisfies $a>19/4$ and the retained validity
conditions hold; in particular every stipulated $a\ge6$ suffices.
\end{corollary}
\begin{proof}
There are at most $C_1(k+1)^3$ clouds at epoch $k$.
The schedule gives $\nu_k^{-3/4}\le C(k+1)^{3/4}$ in the small-error
range, and
$\log(C/(\nu_k\zeta_{k,j}))\le C\log(C(k+1)/\delta)$.
Thus the epoch budget is
$C(k+1)^{15/4}\log(C(k+1)/\delta)$, including the smaller expected
primary-rejection cost. Summation gives \eqref{eq:comp-work} and the
same rejection event as before. No probability argument in
Section~\ref{sec:sequential} depends on the internal smoothing choice:
it uses only genuine uniform geometric enclosures, conditional failure
allowances and deterministically capped call counts, all retained here.

The event that epoch $k$ is executed is measurable before its fresh
randomness. Multiplying the conditional expected epoch cost by the tail
bound at $k-1$ gives a convergent series because
$15/4-a<-1$. The exponentially decaying witness term is summable as well.
The finitely many early epochs contribute a finite amount. Localization
work, apparatus travel and geometric postprocessing remain distinct from
the launch count; no polynomial bit-complexity or minimax claim follows.
\end{proof}

The order reduction is a consequence of the already stated smoothness,
not a stronger identifiability assumption. In particular the original
third-derivative confidence construction and all of its bounds remain
valid. The sharper implementation records explicitly where the six
bounded derivatives improve acquisition rather than being used only for
the regularity of the endpoint density.
